\subsection{Management und Führung}
\label{subsec:ManagementUndFührung}

\par Die Steuerung des Workstreams basierte auf einer kompetenzbasierten Führung und der Etablierung einer belastbaren Projekt-Governance. Aufgrund der Komplexität der Liefergegenstände verfügte ich über weitreichende Befugnisse zur Repriorisierung von Aufgaben und zur Einleitung von Stabilisierungsmaßnahmen bei Lieferverzug. Ein wesentlicher Fokus meiner Führungstätigkeit lag auf dem Skill-Management: Ich validierte kontinuierlich die Kompetenzen der Teilprojektleitenden und förderte durch gezieltes Coaching in Bereichen wie Anforderungsmanagement und Architekturverständnis ein einheitliches Qualitätsniveau. Zur strukturellen Absicherung des Projekterfolgs führte ich ein übergreifendes Governance-Modell für \gls{ricefw} und \gls{aif} ein, welches klare Rollen, Verantwortlichkeiten und Qualitätsstandards für alle internen und externen Beteiligten definierte. Diese Maßnahmen stellten sicher, dass der Workstream trotz der heterogenen Teamzusammensetzung agil und lieferfähig blieb.


\begin{comment}
\par Die Besetzung meiner Rolle als leitender \ac{wsl} im Programm \ac{r2h} erfolgte zum 01.01.2025 durch die Programmleitung aufgrund meiner fachlichen und organisatorischen Eignung. Die Komplexität des \ac{ricefw}-Liefergegenstands, die Vielzahl heterogener Zulieferunternehmen und die technischen Abhängigkeiten erforderten eine kompetenzbasierte Besetzung. Auf persönliche Anforderung der technischen Programmleitung wurde ich als \ac{pm} von der \ac{ves} abgeworben. Zudem war ich federführend in die Evaluierung weiterer Rollen eingebunden und beriet Management und IT hinsichtlich der optimalen Teamzusammensetzung.

\par In Eskalationssituationen verfügte ich über weitreichende Befugnisse: Priorisierung und Repriorisierung von Aufgaben, Weisungsbefugnis gegenüber allen internen und externen Umsetzenden sowie das Recht, Stabilitätsmaßnahmen unmittelbar einzuleiten. Dazu gehörten die Einforderung verbindlicher Zeitpläne, Sonderanalysen, Eskalationen an Programmleitung und Lieferantenmanagement sowie die temporäre Umverteilung von Kapazitäten. Bei kritischen Entscheidungen, wie der Abkehr von der ursprünglich vorgesehenen \ac{aif}-Strategie, verantwortete ich Abstimmung, Entscheidungsvorbereitung und Umsetzung eines alternativen technischen Vorgehens.

\par Als Workstream Lead übernahm ich eine umfassende Management- und Führungsfunktion. Intern verantwortete ich das Reporting in die Programmleitung, in den technischen Steuerkreis, das wöchentliche Transformations-Statusmeeting sowie den Führungskreis der \ac{ais}. Extern vertrat ich den Workstream gegenüber Subunternehmern und Implementierungspartnern, steuerte Liefergegenstände, Qualität und Kapazitäten und stellte die Einhaltung vertraglicher Verpflichtungen sicher. Im Alltag führte ich ein multifirmales Team lateral, stellte Struktur, Transparenz und Zielklarheit her und koordinierte Lieferkomponenten über sieben Unternehmen hinweg.

\par Hinsichtlich des Projektscopes hatte ich maßgeblichen Einfluss auf die Ausgestaltung des \ac{ricefw}-Umfangs und programmweite Entscheidungen. Ich bewertete und priorisierte über 400 Objekte, bereitete Scope-Änderungen für Managemententscheidungen vor und analysierte Auswirkungen neuer Anforderungen wie zusätzlicher SAP-Module (\ac{aif}, \ac{mbc}, \ac{drc}) sowie den Integrationszeitpunkt der Non-SAP-Gesellschaft \ac{adn}. Ebenso begleitete ich die strategische Anpassung des Wellenmodells und bewertete Abhängigkeiten, Risiken und technische Auswirkungen für die Programmleitung.

\par Zur Sicherstellung der Leistungsfähigkeit validierte und verbesserte ich kontinuierlich die Skills der beteiligten Projekt- und Teilprojektleitenden. Dies erfolgte durch gezielte Beobachtungen, Leitung technischer und fachlicher Abstimmungen sowie strukturiertes Coaching in Bereichen wie Anforderungsmanagement, Schnittstellentechnik, Testfallqualität und Architekturverständnis. Gemeinsame Reviews und Wissensformate führten zu einem konsistenten Kompetenzniveau im Workstream.

\par Auf Portfolio-Ebene trug ich durch organisatorische Kompetenz zur Weiterentwicklung projekt- und programmleitender Personen bei. Dazu gehörten die Einführung eines einheitlichen Governance-Modells für \ac{ricefw} und \ac{aif}, die Definition klarer Rollen und Verantwortlichkeiten sowie die Etablierung eines verbindlichen Standards für Liefergegenstände und Qualität. Darüber hinaus beriet ich angrenzende Programme und Projekte zu Abhängigkeitsmanagement, Architekturvorgehen, Testmethodik und Releaseplanung und unterstützte so die Weiterentwicklung der organisatorischen Fähigkeiten im gesamten Portfolio.

%% Second Draft

\par Die Besetzung meiner Rolle als leitender \ac{wsl} im Programm \ac{r2h} erfolgte nach Abschluss der initialen Planungsphase zum 01.01.2025 durch die Programmleitung auf Basis meiner fachlichen und organisatorischen Eignung. Aufgrund der Komplexität des \ac{ricefw}-Liefergegenstands, der Vielzahl heterogener Zulieferunternehmen sowie der technischen Abhängigkeiten in der Zielarchitektur war eine nachweisbar kompetenzbasierte Besetzung erforderlich. Ich wurde explizit und auf persönliche Anforderung der technischen Programmleitung als \ac{pm} von der \ac{ves} abgeworben. Zudem war ich in die Evaluierung weiterer Rollen im Workstream federführend eingebunden und beriet Programm- und IT-Management hinsichtlich der optimalen Zusammensetzung des Teams, sodass sowohl interne als auch externe Mitarbeitende entsprechend ihrer Stärken eingebracht werden konnten.

\par In Eskalationssituationen verfügte ich über weitreichende Befugnisse. Diese umfassten insbesondere die Priorisierung und Repriorisierung von Aufgaben, die Weisungsbefugnis gegenüber allen internen und externen Umsetzenden sowie das Recht, Maßnahmen zur Stabilisierung des Workstreams unmittelbar einzuleiten. Dazu gehörten unter anderem die Einforderung verbindlicher Zeitpläne, das Initiieren von Sonderanalysen, die Eskalation an Programmleitung und Lieferantenmanagement sowie die temporäre Umverteilung von Kapazitäten. In mehreren kritischen Situationen, wie beispielsweise der notwendigen Abkehr von der ursprünglich vorgesehenen \ac{aif}-Strategie, war ich verantwortlich für die Abstimmung, Entscheidungsvorbereitung und Umsetzung eines alternativen technischen Vorgehens.

\par In meiner Rolle nahm ich eine umfassende Management- und Führungsfunktion ein, die sowohl interne als auch externe Vertretung des Workstreams einschloss. Innerhalb des Unternehmens verantwortete ich das Reporting in die Programmleitung, den technischen Steuerkreis, das wöchentliche Transformations-Statusmeeting sowie den internen Führungskreis der \ac{ais}. Gegenüber externen Subunternehmern sowie Implementierungspartnern übernahm ich Verhandlungs- und Steuerungsaufgaben, insbesondere im Hinblick auf Liefergegenstände, Qualität, Kapazitätssteuerung und die Einhaltung vertraglicher Verpflichtungen. Im täglichen Betrieb führte ich ein multi-firmales Team lateral, stellte Struktur, Transparenz und Zielklarheit sicher und sorgte für die zielgerichtete Koordination der Lieferkomponenten über sieben beteiligte Unternehmen hinweg.

\par Im Hinblick auf den Projektscope hatte ich maßgeblichen Einfluss sowohl auf die Ausgestaltung des \ac{ricefw}-Umfangs als auch auf programmweite Entscheidungen. Ich bewertete und priorisierte über 400 geplante Objekte, bereitete Scope-Änderungen für Managemententscheidungen vor und analysierte die Auswirkungen neuer Anforderungen wie der Einführung zusätzlicher SAP-Module (\ac{aif}, \ac{mbc}, \ac{drc}) oder den geeigneten Integrationsansatz und Zeitpunkt der Non-SAP-Gesellschaft \ac{adn}. Ebenso begleitete ich die programmstrategische Anpassung des Wellenmodells (Aufspaltung der Welle 1 ) und brachte meine Einschätzung zu Abhängigkeiten, Risiken und technischen Auswirkungen in die Entscheidungsfindung der Programmleitung ein.

\par Zur Sicherstellung der Leistungsfähigkeit des Workstreams validierte und verbesserte ich kontinuierlich die Skills der beteiligten Projekt- und Teilprojektleitenden. Dies erfolgte durch gezielte Beobachtung der Arbeitsweise, durch Leitung von technischen und fachlichen Abstimmungen sowie durch strukturiertes Coaching, unter anderem in den Bereichen Anforderungsmanagement, Schnittstellentechnik, Testfallqualität und Architekturverständnis. Durch gemeinsame Reviews, Peer-Sessions und abgestimmte Wissensformate wurde ein konsistentes Kompetenzniveau im Workstream erreicht, das zur erfolgreichen Umsetzung der Liefergegenstände beitrug.

\par Auf Portfolio-Ebene trug ich durch organisatorische Kompetenz dazu bei, die Fähigkeiten projekt- und programmleitender Personen zu validieren und zu verbessern. Dies erfolgte insbesondere durch die Einführung eines einheitlichen Governance-Modells für \ac{ricefw} sowie \ac{aif}, die Definition klarer Rollen und Verantwortlichkeiten sowie die Etablierung eines verbindlichen Liefergegenstands- und Qualitätsstandards. Darüber hinaus beriet ich angrenzende Programme und Projekte hinsichtlich Abhängigkeitsmanagement, Architekturvorgehen, Testmethodik und Releaseplanung und förderte so die kontinuierliche Weiterentwicklung der organisationellen Fähigkeiten im gesamten Portfolio.
\end{comment}